\section{Linguistic examples}

Linguistic Examples:

The first sentence is taken from a review and part of the treebank released by \cite{silveira-etal-2014-gold}. 
The second sentence is made-up. 

\begin{exe}
    \ex \label{ex:intro} \begin{xlist} 
        \ex \label{ex:intro-orig} The service was friendly and fast, but this just doesn't make up for the lackluster product.
        \ex \label{ex:intro-madeup} The popcorn was fresh and salty, but this just doesn't make up for the bad movie.
        \ex \label{ex:intro-nonce} The interior was outdated and fresh, but this rather doesn't make up for the newborn period.
    \end{xlist}
\end{exe}

Glossing is also possible and should follow the Leipzig Glossing rules as closely as possible.

\begin{itemize}
    \item \url{https://www.eva.mpg.de/lingua/pdf/Glossing-Rules.pdf}
    \item \url{https://ctan.org/pkg/gb4e}
\end{itemize}

An example with some light glossing: 

\begin{exe}
    \ex \label{ex:man-sees-actress} \begin{xlist}
        \ex \label{ex:man-see-actress-a} \gll Der Mann sieht die Schauspielerin im Fernsehen.\\
        DEF.NOM.SG man see.PRES.3SG DEF actress in TV\\
        `The man sees the actress on TV.'
        \ex \label{ex:man-see-actress-b} \gll Die Schauspielerin sieht der Mann im Fernsehen.\\
        DEF actress see.PRES.3SG DEF.NOM.SG man in TV\\
        `The man sees \textit{the actress} on TV.'
        \ex \label{ex:man-see-actress-c} Die Schauspielerin sieht den Mann im Fernsehen.\\
        `The actress sees the man on TV.'
    \end{xlist}
\end{exe}


\section{Dependency trees}

Using \texttt{tikz-dependency}, example in Figure~\ref{fig:nonce-data-general-algo}. Code for bigger figures and tables should be put in separate files. 

\begin{figure}[ht]
    \centering
    \begin{subfigure}[]{\textwidth}
        \begin{enumerate}
            \item Define the set of UPOS tags that are considered content words. 
            \item For each content word token $t$ in the input treebank, collect the syntactic context $s(t)$. Fill the map $\mathcal{C}$ that maps syntactic contexts to the set of word forms that appear in that context.
            \item Create the output treebank. Replace each content word token $t$ in the input treebank with a random token from the set $\mathcal{C}(s(t))$.
        \end{enumerate}
        \caption{The algorithm}
    \end{subfigure}
    \begin{subfigure}[]{\textwidth}\centering\scalebox{.85}{
        %\includegraphics[width=\textwidth]{img/ptb-noncegeneration-example.png}
        % The service was friendly and fast. 
        % The interior was outdated and fresh
        \begin{tikzpicture}
        \node at (0cm, 0cm) {
        \begin{dependency}
        \begin{deptext}
            \texttt{DET} \& \texttt{NOUN} \& \texttt{AUX} \& \texttt{ADJ} \& \texttt{CCONJ} \& \texttt{ADJ} \& \texttt{PUNCT} \\
            The \& service \& was \& friendly \& and \& fast \& . \\
            The \& \tbf{interior} \& was \& \tbf{outdated} \& and \& \tbf{fresh} \& . \\
        \end{deptext}
        \depedge{2}{1}{det}
        \deproot{4}{root}
        \depedge{4}{2}{nsubj}
        \depedge{4}{3}{cop}
        \depedge{4}{6}{cconj}
        \depedge{4}{7}{punct}
        \depedge{6}{5}{cc}
    \end{dependency}};

    \node at (-40mm, 40mm) {
        \begin{dependency}
            \begin{deptext}
                \& NOUN \\
                \& interior \\
            \end{deptext}
            \deproot{2}{nsubj}
            \depedge{2}{1}{det}
        \end{dependency}};
    
    \node at (-8mm, 40mm) {
        \begin{dependency}
            \begin{deptext}
                \& \& ADJ \& \& \\
                \& \& outdated \& \& \\
            \end{deptext}
            \deproot{3}{root}
            \depedge{3}{1}{nsubj}
            \depedge{3}{2}{cop}
            \depedge{3}{4}{cconcj}
            \depedge{3}{5}{punct}
        \end{dependency}};
    \node at (50mm, 40mm) {
        \begin{dependency}
            \begin{deptext}
                \& ADJ \\
                \& fresh \\
            \end{deptext}
            \deproot{2}{cconj}
            \depedge{2}{1}{cc}
        \end{dependency}};

    \draw[->, dashed] (-38mm, 25mm) -- (-29mm, 0mm);
    \draw[->, dashed] (-10mm, 25mm) -- (-3mm, 9mm);
    \draw[->, dashed] (48mm, 25mm) -- (30mm, 0mm);
\end{tikzpicture}
    }
        \caption{A graphical example}\label{fig:nonce-data-general-algo-example}
    \end{subfigure}
    \caption[Language-independent nonce data generation]{\label{fig:nonce-data-general-algo}Language-independent nonce data generation algorithm. 
    The syntactic context $s(t)$ of the replacements is illustrated by the tree fragments. Furthermore, $\mathcal{C}$ contains the information that $s(service)=s(interior), s(friendly)=s(outdated), s(fast)=s(fresh)$.}
\end{figure}



\section{Tables}

Example table is displayed in Table~\ref{tab:exp1-data}. 
For complex tables, multicolumn and multirow are your friends!

\begin{table}\centering\scalebox{.8}{
\begin{tabular}{l|rrrrrrrrrrrrr}
    \toprule
    {} & train & test & \multicolumn{2}{c}{\texttt{NOUN}} &  \multicolumn{2}{c}{\texttt{PROPN}} & \multicolumn{2}{c}{\texttt{ADJ}} &  \multicolumn{2}{c}{\texttt{ADV}} &  \multicolumn{2}{c}{\texttt{VERB}} & total \\
    & tokens & tokens & freq. & repl. & freq. & repl. & freq. & repl. & freq. & repl. & freq. & repl. & repl. \\
    \midrule
    ar &        224K &        28K &            0.34 &           0.93 &             0.01 &            0.77 &           0.10 &          0.93 &           0.01 &          0.74 &            0.08 &           0.66 &              0.47 \\
    de &       2754K &       326K &            0.21 &           0.96 &             0.06 &            0.93 &           0.08 &          0.68 &           0.06 &          0.92 &            0.08 &           0.61 &              0.41 \\
    en &        204K &        25K &            0.16 &           0.85 &             0.08 &            0.84 &           0.07 &          0.86 &           0.04 &          0.83 &            0.11 &           0.59 &              0.38 \\
    fr &        355K &        10K &            0.19 &           0.83 &             0.05 &            0.69 &           0.06 &          0.88 &           0.05 &          0.92 &            0.08 &           0.50 &              0.36 \\
    ru &       1206K &       158K &            0.23 &           0.84 &             0.04 &            0.70 &           0.09 &          0.80 &           0.05 &          0.89 &            0.11 &           0.46 &              0.39 \\
    zh &         99K &        12K &            0.28 &           0.88 &             0.08 &            0.92 &           0.02 &          0.80 &           0.05 &          0.94 &            0.15 &           0.65 &              0.48 \\
    \bottomrule
\end{tabular}
}
\caption[Data statistics for Experiment 1]{Data statistics for Experiment 1. For each UPOS tag that is affected by the replacement, freq. denotes the relative frequency of the tag in the testing data, and repl. denotes the relative frequency of forms with that tag that are replaced with other forms. The total column shows the relative frequency of forms that are replaced across all tags. zh = Chinese.}
\label{tab:exp1-data}
\end{table}




